\documentclass[10pt,a4paper]{amsart}
\usepackage[margin=19mm]{geometry}
\usepackage{amsmath,amssymb,mathtools}
\usepackage[scr=rsfs,cal=euler]{mathalfa}
\usepackage{xfrac}
\usepackage[unicode,hidelinks,bookmarks=false]{hyperref}
\newcommand{\ie}{i.e.\ }
\newcommand{\cf}{cf.\ }
\newcommand{\resp}{respectively\ }
\newcommand{\Subsection}{\subsection}
\providecommand{\nobreakdash}{\nobreak\hbox{-}}
\setlength{\emergencystretch}{2em}
\multlinegap0pt
\def\baqn{\begin{eqnarray*}}
\def\eaqn{\end{eqnarray*}}
\theoremstyle{plain}
\newtheorem{theorem}{Theorem}[section]
\newtheorem{lemma}[theorem]{Lemma}
\newtheorem{corollary}[theorem]{Corollary}
\newtheorem{proposition}[theorem]{Proposition}
\newtheorem{fact}[theorem]{Fact}
\newtheorem{assumption}[theorem]{Assumption}
\theoremstyle{definition}
\newtheorem{definition}[theorem]{Definition}
\newtheorem{example}[theorem]{Example}
\newtheorem{remark}[theorem]{Remark}
\title[Splitting Algorithms Using Pairs Monotonicity to Solve the Sum of Two Non-monotone Operators Problem]{Splitting Algorithms Using Pairs Monotonicity to Solve the Sum of Two Non-monotone Operators Problem}
\author{Ba Khiet Le, Zakaria Mazgouri,, Michel Th\'era, Ba Khiet Le, 
~
Zakaria Mazgouri, 
~and~ 
Michel Th\' era}
\date{Source: arXiv:2609.27974v1 (CC BY 4.0)}
\begin{document}
\maketitle

\noindent\emph{Source and licence.} The delivered work is the article shown in the title above, version arXiv:2609.27974v1 (CC BY 4.0), distributed under the Creative Commons Attribution 4.0 International licence, as recorded in the arXiv licence metadata for this exact version and shown on the abstract page saved with this item. The full licence notice, which carries the licence URL and the address of that abstract page, is the bundled file \texttt{licenses/arxiv-2609.27974-CC-BY-4.0.txt}.
\par\smallskip

\noindent\textit{Editorial scope.} Counted unit: the article's Theorem (source label thm1) (source0.tex lines 533--545, source label \texttt{thm1}): ``Suppose that Assumptions 1 and 2 hold. If is -strongly monotone and is -Lipschitz continuous, with 0 2 L 2 , then the sequence generated by Algorithm 1 converges to the unique solution . Moreover: (a)''. It is delivered together with the article's complete original proof (source0.tex lines 547--604, one proof environment adjacent to the statement, 58 lines), copied line for line from the article's own text. Required context retained uncounted: none beyond the counted statement and proof. Editorial formatting only: an AMS \texttt{amsart} preamble that restates the article's own macro definitions and its own theorem declarations, and the theorem environments the retained lines use; the article's own class file is neither shipped nor input; the retained environments are renumbered sequentially in order of appearance, whereas the article prints them with its own numbering; the article's equation numbering is not reproduced and every cross-reference inside the retained text resolves to the wrapper's numbering. No citation occurs inside any retained line, so the wrapper carries no bibliography; All retained bodies are exact copies of the bundled \texttt{sources/211539/source0.tex}, so no omission receipt applies. No asset-inclusion command occurs inside any retained span and the package ships no image asset.


\section{The counted unit: Theorem (source label thm1) and its complete original proof}

\begin{theorem}\label{thm1} 
Suppose that Assumptions 1 and 2 hold. If $(A,v)$ is $\alpha$-strongly monotone and $A$ is $L$-Lipschitz continuous, with
\[
0<\gamma<\frac{2\alpha}{L^2},
\]
then the sequence $(x_k)$ generated by Algorithm~1 converges to the unique solution $x^*$.

Moreover:
\begin{enumerate}
\item[(a)] If $v$ is Lipschitz continuous, then $(v(x_k))$ converges linearly to $v(x^*)$.
\item[(b)] If both $v$ and $v^{-1}$ are Lipschitz continuous, then $(x_k)$ converges linearly to $x^*$.
\end{enumerate}
\end{theorem}

\begin{proof}
Since $(A,v)$ is $\alpha$-strongly monotone, the solution set $S$ has a unique point $x^*$.
{Since $v(x_{k+1})=T_{\gamma B}^v(v(x_k)-\gamma Ax_k)$ and $T_{\gamma B}^v$ is nonexpansive, we have
\baqn
\Vert v(x_{k+1})-v(x^*)\Vert^2&\le&\Vert v(x_{k})-v(x^*)-\gamma (Ax_k-Ax^*)\Vert^2\\
&=&\Vert v(x_{k})-v(x^*)\Vert^2-2\gamma\langle v(x_{k})-v(x^*)), Ax_k-Ax^*\rangle+\gamma^2 \Vert Ax_k-Ax^*\Vert^2\\
&\le&\Vert v(x_{k})-v(x^*)\Vert^2-2\gamma \alpha \Vert x_{k}-x^*\Vert^2+\gamma^2L^2\Vert x_{k}-x^*\Vert^2\\
&\le&\Vert v(x_{k})-v(x^*)\Vert^2-\beta \Vert x_{k}-x^*\Vert^2,
\eaqn
where
\[
\beta:=2\gamma\alpha-\gamma^2L^2>0.
\]
Hence, the sequence $\bigl(\|v(x_k)-v(x^*)\|\bigr)$ is monotonically decreasing and therefore convergent.  Consequently $$\|x_k-x^*\|\to0.$$
}
%\medskip

\noindent
 {
(a) Suppose that $v$ is Lipschitz continuous with modulus $L_v$. Then
\[
\|v(x_k)-v(x^*)\|
\le
L_v\|x_k-x^*\|.
\]
%which implies
%$$
%\Vert v(x_{k+1})-v(x^*)\Vert^2\le\Vert v(x_{k})-v(x^*)\Vert^2-\frac{\beta}{L_v^2} \Vert v(x_{k})-v(x^*)\Vert^2=(1-\frac{\beta}{L_v^2})\Vert v(x_{k})-v(x^*)\Vert^2.
%$$
Substituting this estimate into the previous inequality gives
\[
\|v(x_{k+1})-v(x^*)\|^2
\le
\left(1-\frac{\beta}{L_v^2}\right)
\|v(x_k)-v(x^*)\|^2.
\]
Therefore, $(v(x_k))$ converges linearly to $v(x^*)$.
}
\medskip

\noindent
{(b) Assume, in addition, that $v^{-1}$ is Lipschitz continuous with modulus $\widetilde L_v$. Then
\[
\|x_{k+1}-x^*\|
\le
\widetilde L_v\|v(x_{k+1})-v(x^*)\|
\le
\widetilde L_v
\sqrt{1-\frac{\beta}{L_v^2}}
\|v(x_k)-v(x^*)\|
\le
L_v\widetilde L_v
\sqrt{1-\frac{\beta}{L_v^2}}
\|x_k-x^*\|.
\]
Hence, $(x_k)$ converges linearly to $x^*$.
}
\end{proof}

\end{document}
